\section{Semantic Actions dan Konstruksi AST}
\label{sec:ch13-semantic}

\subsection{Nilai Semantik}

Bison menyediakan variabel khusus:
\begin{itemize}[leftmargin=*]
  \item \code{\$\$}: nilai hasil nonterminal kiri;
  \item \code{\$1}, \code{\$2}, ...: nilai simbol pada ruas kanan;
  \item \code{\$<token>1}: akses anggota union tertentu.
\end{itemize}

Aksi dijalankan ketika aturan direduksi. Pada titik itu, semua nilai
simbol ruas kanan sudah tersedia, sehingga AST dapat dibangun dari
bawah ke atas.

\subsection{Membangun AST}

Aturan \code{assign\_or\_call} dan \code{block} menunjukkan pola
pembangunan AST. Potongannya ditunjukkan pada
\lstref{lst:ch13-block}.

\includecode[
  language=Bison,
  firstline=135,
  lastline=166,
  firstnumber=135,
  caption={Aksi untuk assignment, call, dan block},
  label={lst:ch13-block}
]{../../code/subsetc/src/parser.y}

Perhatikan bahwa aksi tidak hanya membuat simpul, tetapi juga
memindahkan anak dari daftar \code{next} ke daftar \code{children}.
Pemisahan ini menjaga agar AST memiliki struktur pohon yang bersih.

\subsection{Contoh AST}

Aturan \code{primary} membangun simpul literal, identifier, dan call.
Potongannya ditunjukkan pada \lstref{lst:ch13-primary}.

\includecode[
  language=Bison,
  firstline=281,
  lastline=299,
  firstnumber=281,
  caption={Aksi untuk \code{primary}},
  label={lst:ch13-primary}
]{../../code/subsetc/src/parser.y}

Untuk ekspresi \(a + b * c\), aksi \code{term} membangun simpul
\code{*} dengan anak \code{b} dan \code{c}; aksi \code{additive}
membangun simpul \code{+} dengan anak \code{a} dan simpul \code{*}.
Hasilnya adalah AST yang mencerminkan precedence.

\begin{figure}[H]
\centering
\begin{forest}
asttree
[+
  [\code{a}]
  [*
    [\code{b}]
    [\code{c}]
  ]
]
\end{forest}
\caption{AST untuk \(a + b * c\)}
\label{fig:ch13-ast}
\end{figure}

\subsection{Location Tracking}

Bison dapat melacak lokasi token dan aturan. Dengan
\code{\%locations}, setiap simbol memiliki \code{yylloc} yang dapat
disalin ke simpul AST. Lokasi ini berguna untuk diagnostik semantik
pada Bab XIV--XVI.

\subsection{Pesan Kesalahan Sintaksis}

Fungsi \code{yyerror} dipanggil Bison ketika terjadi kesalahan.
Pesan yang baik memuat lokasi, penyebab, dan token yang ditemukan.
Dengan \code{\%locations}, \code{yylloc} menyediakan baris dan kolom
yang dapat dipakai langsung.

\begin{lstlisting}[language=C]
void yyerror(const char *msg) {
    fprintf(stderr, "%d:%d: %s\n",
            yylloc.first_line, yylloc.first_column, msg);
}
\end{lstlisting}

\begin{table}[H]
\centering
\caption{Komponen pesan kesalahan sintaksis}
\label{tab:ch13-diagnostik}
\begin{tabularx}{\textwidth}{@{}l X@{}}
\toprule
\textbf{Komponen} & \textbf{Contoh} \\
\midrule
Lokasi & baris 7, kolom 5 \\
Penyebab & diharapkan \code{;} setelah ekspresi \\
Token ditemukan & \code{\}} \\
Saran perbaikan & tambahkan \code{;} sebelum \code{\}} \\
\bottomrule
\end{tabularx}
\end{table}

Pesan ini diuji pada kasus valid dan tidak valid. Pesan yang hanya
berisi ``syntax error'' tidak memenuhi Sub-CPMK 4.4 karena tidak
membantu pengguna menemukan kesalahan.

\begin{infobox}
Semantic action sebaiknya singkat dan fokus pada pembangunan AST.
Logika kompleks sebaiknya dipindahkan ke fungsi bantu agar grammar
tetap mudah dibaca.
\end{infobox}
